%% parts/p06/constructibility-forcing.tex --- constructibility, forcing,
%% and exact relative-consistency status. Source: v1.1
%% ch:constructibility-forcing. Body statements only.

\chapter{Constructibility, Forcing, and Independence}
\label{sec:p06-constructibility-forcing}
\chaptermark{Constructibility and Forcing}

This chapter keeps two routes sharply distinct.  The constructible universe
supplies the positive relative-consistency result locally.  The negative-CH
result remains an imported Boolean-valued theorem.  Countable-transitive-model
forcing is retained only as a pedagogical method under its explicitly stronger
hypothesis.  The standard source is \cite{Jech2003}.

%% parts/p06/forcing-vocabulary.tex --- shared forcing vocabulary before use

\section{Forcing-order vocabulary}
\label{sec:p06-forcing-vocabulary}

The exact source convention at
Definition~\ref{def:p06-forcing-order-orientation} fixes the
stronger-is-smaller orientation, and the exact source definition at
Definition~\ref{def:p06-poset-filter} controls poset filters, dense sets, and
genericity.  This lookup does not create second identities for them.  It adds
only the cardinal notation needed before Chapter~6.5, compatibility/ccc, and
the many-real Cohen poset used by the later supplements
\cite[pp.~201--224]{Jech2003}.

\begin{hbconvention}{def:p06-forcing-cardinal-notation}{Cardinal notation used before Chapter 6.5}
The target architecture places forcing before the full cardinal-arithmetic
entry.  To prevent a forward use from being undefined, only the notation
needed in Chapters~6.3--6.4 is fixed here.  For cardinals
\(\kappa,\lambda\),
\[
  \kappa^\lambda\defeq\abs{{}^{\lambda}\kappa},
  \qquad
  \lambda^{<\kappa}\defeq
    \sup_{\mu<\kappa}\lambda^\mu,
\]
where \({}^{\lambda}\kappa\) is the set of functions from \(\lambda\) to
\(\kappa\), and \(\kappa^+\) is the least cardinal strictly greater than
\(\kappa\).  The alephs needed below are fixed by
\[
  \aleph_0\defeq\omega,
  \qquad
  \aleph_{\alpha+1}\defeq(\aleph_\alpha)^+,
  \qquad
  \aleph_\delta\defeq\sup_{\beta<\delta}\aleph_\beta
  \quad(\delta\text{ a nonzero limit ordinal}).
\]
Thus
\[
  \mathsf{CH}\defeq2^{\aleph_0}=\aleph_1,
  \qquad
  \mathsf{GCH}\defeq
    \forall\alpha\,\bigl(2^{\aleph_\alpha}=\aleph_{\alpha+1}\bigr).
\]
For an ordinal \(\alpha\), a set \(C\subseteq\alpha\) is
\emph{cofinal} iff every \(\beta<\alpha\) is at most some member of \(C\),
and \(\operatorname{cf}(\alpha)\) is the least order type of such a set.
An infinite cardinal \(\kappa\) is \emph{regular} iff
\(\operatorname{cf}(\kappa)=\kappa\).  These are imported lookup clauses,
not theorem or status claims; Definition~\ref{def:p06-cardinal-arithmetic}
and Definition~\ref{def:p06-cofinality-regular-singular} give the full
Chapter~6.5 entries.
\end{hbconvention}

\begin{hbdefinition}{def:p06-ccc}{Compatibility, antichains, and ccc}
Conditions \(p,q\in\mathbb P\) are \emph{compatible} iff some
\(p^\ast\in\mathbb P\) satisfies \(p^\ast\leq p\) and \(p^\ast\leq q\).  An
\emph{antichain} is a subset of pairwise incompatible conditions.  The poset
has the \emph{countable chain condition} (ccc) iff every antichain is
countable.  These are metalinguistic properties of a poset, not filter clauses;
the distinction is used by the preservation and name-counting theorems.
\index{countable chain condition}\index{antichain!forcing}
\end{hbdefinition}

\begin{hbdefinition}{def:p06-add-forcing}{Cohen forcing \(\operatorname{Add}(\omega,\kappa)\)}
For a cardinal \(\kappa\),
\[
  \operatorname{Add}(\omega,\kappa)
  \defeq
  \operatorname{Fn}(\kappa\times\omega,2,{<}\omega),
\]
the finite partial functions from \(\kappa\times\omega\) to \(2\), ordered by
reverse inclusion.  This poset adds \(\kappa\) coordinate Cohen reals; it is
not the one-real forcing unless \(\kappa=1\).  The distinction is needed
because adding one Cohen real does not by itself force \(\neg\mathsf{CH}\).
\index{Cohen forcing}\index{Add forcing@\(\operatorname{Add}(\omega,\kappa)\)}
\end{hbdefinition}


\section{The constructible hierarchy}
\label{sec:p06-constructible-hierarchy}

\begin{hbdefinition}{def:p06-definable-subsets}{Definable subsets of a set}
For a set \(A\), let
\[
  \operatorname{Def}(A)
  \defeq
  \left\{X\subseteq A:
    \begin{array}{l}
      \text{for some first-order formula }\varphi(y,\vec z)
      \text{ and }\vec a\in A^{<\omega},\\
      X=\{x\in A:(A,\in),s_{x,\vec a}\entails\varphi\}
    \end{array}
  \right\}.
\]
Here \(s_{x,\vec a}\) assigns \(x\) to \(y\) and the entries of
\(\vec a\) to the displayed parameter variables.  Satisfaction is the
ordinary set-sized relation; this definition introduces no truth predicate
for the universe.
\index{definable subset}\index{Def operation@\(\operatorname{Def}\)}
\end{hbdefinition}

\begin{hbdefinition}{def:p06-constructible-hierarchy}{Constructible hierarchy}
Transfinite recursion defines
\[
  L_0=\varnothing,\qquad
  L_{\alpha+1}=\operatorname{Def}(L_\alpha),\qquad
  L_\lambda=\bigcup_{\beta<\lambda}L_\beta
\]
for limit \(\lambda\).  Put
\[
  L=\bigcup_{\alpha\in\operatorname{Ord}}L_\alpha.
\]
A set is \emph{constructible} iff it belongs to \(L\).
\index{constructible hierarchy}\index{constructible set}
\end{hbdefinition}

\begin{hbconvention}{def:p06-constructible-relativization}{Relativization to \(L\)}
For each standard membership-language formula \(\varphi\), define
\(\varphi^L\) metalinguistically: atomic formulas are unchanged, the
operation commutes with the propositional connectives, and
\[
  (\exists x\,\psi)^L\defeq
    \exists x\,(x\in L\land\psi^L),
  \qquad
  (\forall x\,\psi)^L\defeq
    \forall x\,(x\in L\to\psi^L).
\]
The predicate \(x\in L\) is eliminated by the formula defining
constructibility.  Thus \(L\entails\sigma\) abbreviates the ambient assertion
\(\sigma^L\), formula by formula; it neither treats \(L\) as a set-sized
structure nor posits a class truth predicate.
\end{hbconvention}

\begin{hbtheorem}{thm:p06-constructible-class}{sketched=annex:p6.force.1}{The constructible class}
Assume \(\mathsf{ZF}\) in the ambient universe.  The class \(L\) is transitive
and contains every ordinal.
\end{hbtheorem}

\begin{hbtheorem}{thm:p06-l-model}{cited=Jech2003}{Constructibility theorem scheme}
For every standard axiom or schema instance \(\sigma\) of
\(\mathsf{ZFC}+\mathsf{GCH}+(V=L)\),
\[
  \mathsf{ZF}\derives_{\mathsf{ND}_{\mathsf{FOL}}}\sigma^L.
\]
This formula-by-formula scheme includes separation and replacement and does
not assert a single internal satisfaction predicate.
\end{hbtheorem}

\begin{hbtheorem}{thm:p06-condensation}{cited=Jech2003}{Condensation for \(L\)}
Let \(\alpha>\omega\) be a limit ordinal and let
  \(X\preccurlyeq(L_\alpha,\in)\) be a set-sized elementary substructure.  The
Mostowski collapse of \((X,\in)\) is \((L_\beta,\in)\) for some
\(\beta\leq\alpha\).  Elementarity is load-bearing: arbitrary subsets of
\(L_\alpha\) do not satisfy this conclusion.
\end{hbtheorem}

Jensen's fine-structure analysis is a further-reading pointer rather than a
local entry, as required by exclusion row N-x-004.  The reflection component
of row N-a-067 is the finite-family L\'evy reflection statement in
Theorem~\ref{thm:p06-levy-reflection}; its formula-family and parameter
hypotheses are stated there rather than weakened here.

\begin{figure}[htbp]
\centering
\begin{tikzpicture}[scale=.85,transform shape,>=Stealth,
  level/.style={draw,rounded corners,minimum width=3.6cm,
    minimum height=.62cm,fill=blue!5}]
  \node[level] (l0) at (0,0) {$L_0=\varnothing$};
  \node[level] (l1) at (0,1.05)
    {$L_{\alpha+1}=\operatorname{Def}(L_\alpha)$};
  \node[level] (ll) at (0,2.10)
    {$L_\lambda=\bigcup_{\beta<\lambda}L_\beta$};
  \node[level,minimum width=5.3cm] (linf) at (0,3.15)
    {$L=\bigcup_{\alpha\in\operatorname{Ord}}L_\alpha\subseteq V$};
  \draw[->] (l0) -- (l1);
  \draw[->] (l1) -- (ll);
  \draw[->] (ll) -- (linf);
  \node[align=left,anchor=west] at (3.1,1.58)
    {successor stages retain only\\subsets definable over the prior level};
  \draw[->,dashed] (2.95,1.58) -- (1.85,1.15);
  \node[align=center] at (0,4.0)
    {Condensation identifies suitable elementary substructures\\
     with earlier levels; \(V=L\) asserts \(V=L\).};
\end{tikzpicture}
\caption{The constructible hierarchy inside the cumulative universe (A-131):
each \(L_\alpha\subseteq V_\alpha\), and the union defining \(L\) continues
through every ordinal, so \(L\) is a proper class.}
\label{fig:p06-constructible-hierarchy}
\end{figure}

\section{The countable-transitive-model method}
\label{sec:p06-ctm-forcing}

\begin{hbconvention}{def:p06-forcing-standing}{Forcing convention}
Fix externally a countable transitive set \(M\) such that
\((M,\in)\entails\mathsf{ZFC}\), and fix in \(M\) a poset
\((\mathbb P,\leq,1_{\mathbb P})\) with the orientation of
Convention~\ref{def:p06-forcing-order-orientation}.  A poset filter
\(G\subseteq\mathbb P\) is \emph{\(M\)-generic} iff it meets every dense
subset of \(\mathbb P\) that belongs to \(M\).  This exact qualifier is needed
because no filter generic over the ambient universe is being asserted.

The pedagogical hypothesis on \(M\) is stronger than
the arithmetical sentence \(\mathsf{Con}(\mathsf{ZFC})\); it is not assumed in
the final relative-consistency argument for \(\neg\mathsf{CH}\).
\end{hbconvention}

Write \(\mathsf{Con}(\mathsf{ZFC})^\in\) for the membership-language sentence
obtained by translating the fixed arithmetical consistency sentence along the
canonical interpretation of arithmetic by the finite ordinals.  The
superscript distinguishes the interpreted membership-language sentence from
the arithmetical formula itself; the notation is used only in the following
strength comparison.

\begin{hbtheorem}{thm:p06-ctm-strength}{sketched=annex:p6.force.2}{Strength of the countable-transitive-model hypothesis}
Under the canonical interpretation of arithmetic by the finite ordinals, the
existence of this \(M\) implies consistency of
\(\mathsf{ZFC}+\mathsf{Con}(\mathsf{ZFC})^\in\).  Consequently the fixed
hypothesis strictly exceeds \(\mathsf{Con}(\mathsf{ZFC})\), assuming the
standard transfer of the derivability conditions across that interpretation.
\end{hbtheorem}

\begin{hblemma}{thm:p06-bounded-absoluteness}{sketched=annex:p6.force.3}{Bounded absoluteness for the fixed transitive set}
If \(\varphi(\vec x)\) is a membership-language formula whose quantifiers are
all bounded and \(\vec a\in M\), then
\[
  (M,\in)\entails\varphi(\vec a)
  \quad\Leftrightarrow\quad
  \varphi(\vec a)
\]
in the ambient universe.  This is the only transfer between \(M\) and the
ambient universe used by the elementary forcing construction.
\end{hblemma}

\begin{hbmethod}{def:p06-generic-filter-method}{Generic-filter construction}
Externally enumerate the dense subsets of \(\mathbb P\) belonging to the
countable set \(M\).  Starting below a prescribed condition, choose a
descending sequence meeting them in turn; its upward closure is directed and
\(M\)-generic.  This construction is available only under
Convention~\ref{def:p06-forcing-standing}.
\end{hbmethod}

\begin{hblemma}{thm:p06-generic-filter-existence}{proved=annex:p6.force.4}{Existence of a generic filter under the fixed hypothesis}
For every \(p\in\mathbb P\), an \(M\)-generic poset filter containing \(p\)
exists in the ambient universe.
\end{hblemma}

\begin{hbdefinition}{def:p06-names-extension}{Names, evaluation, and extension}
The hierarchy of \(\mathbb P\)-names is
\[
  V^{\mathbb P}_\alpha
  \defeq
  \mathcal P\!\left(
    \left(\bigcup_{\beta<\alpha}V^{\mathbb P}_\beta\right)
    \times\mathbb P
  \right),
  \qquad
  V^{\mathbb P}\defeq
  \bigcup_{\alpha\in\operatorname{Ord}}V^{\mathbb P}_\alpha.
\]
For a poset filter \(G\), rank recursion gives
\[
  \nu^G\defeq
  \{\zeta^G:\exists p\in G\,(\langle\zeta,p\rangle\in\nu)\}.
\]
For \(x\in M\), the \emph{check name}
\[
  \check x\defeq
  \{\langle\check y,1_{\mathbb P}\rangle:y\in x\}
\]
is needed because \(x\) need not have the name--condition pair structure.
Its standing is a \(\mathbb P\)-name encoding the fixed ground-model set
\(x\), not \(x\) itself.  This nonredundant distinction embeds \(M\) into the
extension and supplies the canonical names used in the Cohen and nice-name
entries.  The generic extension is
\[
  M[G]\defeq
  \{\nu^G:\nu\in M\cap V^{\mathbb P}\}.
\]
\index{forcing name}\index{check name}\index{generic extension}
\end{hbdefinition}

\begin{hblemma}{thm:p06-check-name-evaluation}{sketched=annex:p6.force.5}{Evaluation of check names}
For every \(x\in M\),
\[
  \check x\in M\cap V^{\mathbb P}
  \qquad\text{and}\qquad
  (\check x)^G=x.
\]
Thus every member of \(M\) has a name whose value belongs to \(M[G]\).
\end{hblemma}

\begin{hbdefinition}{def:p06-semantic-forcing}{Semantic forcing}
Let the free variables of \(\varphi\) lie among pairwise distinct
\(x_1,\ldots,x_n\), let
\(\nu_1,\ldots,\nu_n\in M\cap V^{\mathbb P}\), and let \(p\in\mathbb P\).
Then
\[
  p\forces_{M,\mathbb P}
  \varphi(\nu_1,\ldots,\nu_n)
\]
  iff every \(M\)-generic \(G\ni p\), under every assignment sending \(x_i\) to
\(\nu_i^G\), makes \(\varphi\) true in \((M[G],\in)\).  This is an external
metalinguistic relation, neither ordinary satisfaction nor the recursively
defined forcing relation imported below.  The qualifier is therefore
necessary, and the relation is used in the monotonicity lemma and in the
forcing-theorem comparison.
\index{forcing!semantic}
\end{hbdefinition}

\begin{hblemma}{thm:p06-semantic-forcing-monotone}{proved=annex:p6.force.6}{Monotonicity of semantic forcing}
If \(q\leq p\) and
\(p\forces_{M,\mathbb P}\varphi(\vec\nu)\), then
\(q\forces_{M,\mathbb P}\varphi(\vec\nu)\).
\end{hblemma}

\begin{hbtheorem}{thm:p06-forcing-theorem}{cited=Jech2003}{Forcing theorem under the fixed hypothesis}
For each fixed standard membership-language formula \(\varphi(\vec x)\), a
formula with the fixed poset as parameter represents the recursively defined
forcing assertion over \(M\).  For names \(\vec\nu\in M\) and an
\(M\)-generic \(G\), recursive forcing agrees with semantic forcing; the truth
lemma says
\[
  (M[G],\in)\entails\varphi(\vec\nu^G)
  \quad\Leftrightarrow\quad
  \exists p\in G\;p\forces\varphi(\vec\nu);
\]
and \(M[G]\) satisfies \(\mathsf{ZFC}\) and has the same ordinals as \(M\).
Definability is formula by formula, not a single truth predicate for all
standard formula codes.
\end{hbtheorem}

\begin{proofboundary}{def:p06-forcing-proof-boundary}{Boundary of the forcing citation}
A local proof would have to define the two simultaneous atomic clauses, give
every recursive logical clause and its measure, prove density and truth lemmas
constructor by constructor, and verify every \(\mathsf{ZFC}\) axiom in the
extension, including Replacement.  Those steps are not reproduced; the
recursive forcing theorem is cited and has no annex item.
\end{proofboundary}

\begin{hbconvention}{def:p06-recursive-forcing-use}{Imported recursive forcing}
After Theorem~\ref{thm:p06-forcing-theorem}, an unadorned
\(p\forces\varphi(\vec\nu)\) denotes the cited recursive relation.  The
indexed form \(\forces_{M,\mathbb P}\) continues to denote
Definition~\ref{def:p06-semantic-forcing}.  This abbreviation is needed for
the decision-antichain statements and introduces no second forcing glyph.
\end{hbconvention}

\begin{hbtheorem}{thm:p06-cohen-new-real}{proved=annex:p6.force.7}{One Cohen real}
Let
\(\mathbb P=\operatorname{Fn}(\omega,2,{<}\omega)\)
be ordered by reverse inclusion, and let \(G\) be \(M\)-generic.  Then
\(c=\bigcup G\) is a total function \(\omega\to2\), belongs to \(M[G]\), and
does not belong to \(M\).  This result alone does not imply
\(\neg\mathsf{CH}\).
\end{hbtheorem}

\begin{figure}[htbp]
\centering
\begin{tikzpicture}[scale=.74,transform shape,>=Stealth,node distance=2.4cm,
  box/.style={draw,rounded corners,minimum width=3.0cm,
    minimum height=1.05cm,align=center}]
  \node[box,fill=blue!5] (m) {\(M\)\\names \(\nu\), poset \(\mathbb P\)};
  \node[box,fill=green!7,right=of m] (g)
    {external \(M\)-generic\\filter \(G\subseteq\mathbb P\)\\
     \(G\notin M\) in the Cohen case};
  \node[box,fill=orange!8,right=of g] (mg)
    {\(M[G]\)\\values \(\nu^G\)};
  \node[box,fill=yellow!8,below=1.0cm of g] (ev)
    {name evaluation\\\((\nu,G)\mapsto\nu^G\)};
  \draw[->] (m) -- node[above]{dense sets} (g);
  \draw[->] (m.south east) -- node[below left]{names} (ev.west);
  \draw[->] (g.south) -- node[right]{parameter} (ev.north);
  \draw[->] (ev.east) -- node[below right]{values} (mg.south west);
  \draw[->,bend left=25] (m) to
    node[above]{\(x\mapsto\check x\mapsto x\)} (mg);
  \node[below=1.0cm of ev,align=center]
    {The pedagogical construction assumes external countability.\\
     Cited forcing theorem: \(M[G]\entails\mathsf{ZFC}\) and
     \(\operatorname{Ord}^{M[G]}=\operatorname{Ord}^M\).};
\end{tikzpicture}
\caption{Names, a generic filter, and the forcing extension (A-132).}
\label{fig:p06-forcing-extension}
\end{figure}

\section{Chain condition and many Cohen reals}
\label{sec:p06-ccc-cohen}

\begin{hbtheorem}{thm:p06-ccc-preservation}{proved=annex:p6.force.9}{Cardinal and cofinality preservation by ccc forcing}
Assume \((M,\in)\entails\mathsf{ZFC}\), \(M\) regards \(\mathbb P\) as ccc,
and \(G\) is \(M\)-generic.  Then every \(M\)-cardinal remains a cardinal of
\(M[G]\), and every ordinal has the same cofinality as computed in \(M\) and
\(M[G]\).
\end{hbtheorem}

\begin{hbdefinition}{def:p06-nice-name}{Nice name for a subset of \(\omega\)}
The adjective \emph{nice} isolates the antichain data needed to count names;
an arbitrary name does not display that data.  A \emph{nice
\(\mathbb P\)-name for a subset of \(\omega\)} is
\[
  \dot x=\bigcup_{n<\omega}
  \{\langle\check n,p\rangle:p\in A_n\},
\]
where every \(A_n\) is an antichain in \(\mathbb P\).  This is a
metalinguistic subclass of the names in
Definition~\ref{def:p06-names-extension}; it does not change evaluation.  Its
nonredundant use is that ccc makes each \(A_n\) countable.
\index{nice name!forcing}
\end{hbdefinition}

\begin{hbtheorem}{thm:p06-nice-name-counting}{proved=annex:p6.force.10}{Nice-name reduction and counting}
Assume \(M\) regards \(\mathbb P\) as ccc.  If \(p\in\mathbb P\) and
\(p\forces\nu\subseteq\check\omega\), then some nice
\(\mathbb P\)-name \(\zeta\) satisfies \(p\forces\nu=\zeta\).
Consequently, if \(\mathbb P\) is infinite and
\(\lambda=(\abs{\mathbb P}^{\aleph_0})^M\), then
\[
  M[G]\entails 2^{\aleph_0}\leq\lambda.
\]
\end{hbtheorem}

\begin{hbcorollary}{thm:p06-many-cohen-continuum}{proved=annex:p6.force.11}{Forcing the size of the continuum}
Suppose \(M\) regards \(\kappa\) as an infinite cardinal and satisfies
\(\kappa^{\aleph_0}=\kappa\), and let
\(\mathbb P=\operatorname{Add}(\omega,\kappa)^M\).  For every
\(M\)-generic \(G\),
\[
  M[G]\entails 2^{\aleph_0}=\kappa.
\]
In particular, a ground model of \(\mathsf{GCH}\) with
\(\kappa=\aleph_2\) yields a forcing extension satisfying
\(\neg\mathsf{CH}\).
\end{hbcorollary}

\section{Exact relative-consistency status}
\label{sec:p06-relative-consistency}

\begin{hblemma}{thm:p06-relativization-derivations}{proved=annex:p6.force.12}{Relativization of derivations}
Let \(\varphi_1,\ldots,\varphi_k,\psi\) have free variables among pairwise
distinct \(x_1,\ldots,x_m\).  If
\(\{\varphi_1,\ldots,\varphi_k\}
  \derives_{\mathsf{ND}_{\mathsf{FOL}}}\psi\) in
\hyperref[sheet:nd-fol]{\texttt{sheet:nd-fol}}, then
\[
  \mathsf{ZF}\derives_{\mathsf{ND}_{\mathsf{FOL}}}
  \forall x_1\cdots\forall x_m
  \left(
    \bigwedge_{1\leq j\leq m}x_j\in L\land
    \bigwedge_{1\leq i\leq k}\varphi_i^L
    \to\psi^L
  \right),
\]
with the corresponding conjunction absent when \(m=0\) or \(k=0\).
\end{hblemma}

\begin{hbtheorem}{thm:p06-ch-relative-consistency}{proved=annex:p6.force.13}{Positive relative consistency of CH}
If \(\mathsf{ZF}\) is consistent, then
\(\mathsf{ZFC}+\mathsf{GCH}+(V=L)\) is consistent.  Hence, if
\(\mathsf{ZFC}\) is consistent, then
\(\mathsf{ZFC}+\mathsf{CH}\) is consistent.
\end{hbtheorem}

\begin{hbtheorem}{thm:p06-not-ch-relative-consistency}{cited=Jech2003}{Relative consistency of the negation of CH}
If \(\mathsf{ZFC}\) is consistent, then
\(\mathsf{ZFC}+\neg\mathsf{CH}\) is consistent.  The imported proof uses the
Boolean completion of
\(\operatorname{Fn}(\aleph_2\times\omega,2,{<}\omega)\): the Boolean-valued
model satisfies \(\mathsf{ZFC}\), the forcing is ccc and preserves
\(\aleph_1,\aleph_2\), and the \(\aleph_2\)-indexed Cohen sections force
\(2^{\aleph_0}\geq\aleph_2\).  This cited theorem, not the countable
transitive model convention, discharges the consistency implication.
\end{hbtheorem}

\begin{hbtheorem}{thm:p06-arithmetized-ch-consistency}{sketched=annex:p6.force.14}{Arithmetized CH consistency}
For the fixed proof predicate and canonical interpretation of arithmetic in
set theory, \(\mathsf{PA}\) proves
\[
  \mathsf{Con}(\mathsf{ZFC})\to
  \mathsf{Con}(\mathsf{ZFC}+\mathsf{CH})
\]
and
\[
  \mathsf{Con}(\mathsf{ZFC})\to
  \mathsf{Con}(\mathsf{ZFC}+\neg\mathsf{CH}).
\]
\end{hbtheorem}

\begin{hbtheorem}{thm:p06-ch-nonderivability}{proved=annex:p6.force.15}{Nonderivability of CH and its negation}
If \(\mathsf{ZFC}\) is consistent, then
\[
  \mathsf{ZFC}\notderives_{\mathsf{ND}_{\mathsf{FOL}}}\mathsf{CH}
  \qquad\text{and}\qquad
  \mathsf{ZFC}\notderives_{\mathsf{ND}_{\mathsf{FOL}}}\neg\mathsf{CH}.
\]
\end{hbtheorem}

\begin{hbwarning}{def:p06-consistency-boundary}{Consistency limits}
Neither direction proves \(\mathsf{Con}(\mathsf{ZFC})\) without a consistency
hypothesis.  Nor may the construction from a countable transitive model
replace the Boolean-valued discharge: external countability and transitivity
are stronger than the displayed arithmetical consistency sentence.
\end{hbwarning}

\paragraph{Source disposition.}
The two source exercises are deliberately retired because the handbook has no
exercise layer.  Their reusable content---density of the Cohen coordinate
sets, single-valuedness of \(\bigcup G\), and the distinction among
\(L\entails\mathsf{CH}\), formal derivability, and relative
consistency---is stated above.  Product, iterated, and proper forcing remain
forward pointers at the ceilings fixed by D-P7-011 and N-a-068; no
unregistered substantive entry is introduced for HOD.
